% ======================================================================
% Chapter 16: Quantization of Non-Abelian Gauge Theories
% (book pages 505-543; PDF pages 528-567)
% ======================================================================
\chapter{Quantization of Non-Abelian Gauge Theories}

The previous chapter showed how to construct Lagrangians with non-Abelian
gauge symmetry.  However, this is only the first step in the process of
relating the idea of non-Abelian gauge invariance to the real interactions of
particle physics.  We must next work out the rules for computing Feynman
diagrams containing the non-Abelian gauge vector particles, then use these
rules to compute scattering amplitudes and cross sections.  This chapter will
develop the technology needed for such calculations.

Alongside this technical discussion, we will study how the gauge symmetry
affects the Feynman amplitudes.  In any theory with a local symmetry, some
degrees of freedom of the fields that appear in the Lagrangian are
\emph{unphysical}, in the sense that they can be adjusted arbitrarily by gauge
transformations.  In electrodynamics, the components of the field $A^\mu(k)$
proportional to $k^\mu$ lie along the symmetry directions.  We saw in
Section~9.4 that this fact has two important consequences.  First, the
propagator of the field $A^\mu$ is ambiguous; there are multiple expressions
for the propagator, which follow equally well from the QED Lagrangian.
Second, the vertices of electrodynamics are such that this ambiguity makes no
difference in the calculation of cross sections.  For example,
Eq.~\eqref{eq:9.58} displays a continuous family of photon propagators, one
for each value of the continuous parameter $\xi$; but we saw immediately that
all dependence of $S$-matrix elements on $\xi$ is eliminated by the Ward
identity.  Non-Abelian gauge theories contain similar ambiguities and
cancellations, but, as we will see in this chapter, the structure of the
cancellations is more intricate.

An additional goal of this chapter is to compute the Callan-Symanzik
$\beta$ function, and hence determine the behavior of the running coupling
constant, for non-Abelian gauge theories.  As discussed in Chapter~14, these
theories are in fact \emph{asymptotically free}: The coupling constant
becomes weak at large momenta.  This result indicates the applicability of
non-Abelian gauge theory to model the strong interactions.  We will be able
to derive this result once we have determined the correct Feynman rules for
non-Abelian gauge theories.

\section{Interactions of Non-Abelian Gauge Bosons}

Most of the Feynman rules for non-Abelian gauge theory can be read directly
from the Yang-Mills Lagrangian, following the method of Section~9.2.
However, when we quantized the electromagnetic field in Section~9.4, we saw
that the functional integral over a gauge field must be defined carefully, and
that the subtle aspects of this construction can introduce new ingredients
into the quantum theory.  In this section we will see how far we can go in
the non-Abelian theory by ignoring these subtleties.  In Section~16.2 we will
carry out a more proper derivation of the Feynman rules, through a careful
analysis of the functional integral.

\subsection{Feynman Rules for Fermions and Gauge Bosons}

The Yang-Mills Lagrangian, as derived in the previous chapter, is
\begin{equation}
  \mathcal{L} = -\frac{1}{4}(F^a_{\mu\nu})^2 + \bar{\psi}(i\slashed{D}
    - m)\psi ,
  \label{eq:16.1}
\end{equation}
where the index $a$ is summed over the generators of the gauge group $G$, and
the fermion multiplet $\psi$ belongs to an irreducible representation $r$ of
$G$.  The field strength is
\begin{equation}
  F^a_{\mu\nu} = \partial_\mu A^a_\nu - \partial_\nu A^a_\mu
    + gf^{abc}A^b_\mu A^c_\nu ,
  \label{eq:16.2}
\end{equation}
where $f^{abc}$ are the structure constants of $G$.  The covariant derivative
is defined in terms of the representation matrices $t^a$ by
\begin{equation}
  D_\mu = \partial_\mu - igA^a_\mu t^a .
  \label{eq:16.3}
\end{equation}
From now on we will drop the subscript $r$ except where it is needed for
clarity.

The Feynman rules for this Lagrangian can be derived from a functional
integral over the fields $\psi$, $\bar{\psi}$, and $A^a_\mu$.  Imagine
expanding the functional integral in perturbation theory, starting with the
free Lagrangian, at $g = 0$.  The free theory contains a number of free
fermions equal to the dimension $d(r)$ of the representation $r$, and a
number of free vector bosons equal to the number $d(G)$ of generators of
$G$.  Using the methods of Section~9.5, it is straightforward to derive the
fermion propagator
\begin{equation}
  \langle\psi(x)\bar{\psi}(y)\rangle = S_F(x-y)\,\delta^{ij}
  = \int \frac{d^4p}{(2\pi)^4}\ \frac{i(\slashed{p}+m)}{p^2-m^2+i\epsilon}\
    e^{-ip\cdot(x-y)}\,\delta^{ij} ,
  \label{eq:16.4}
\end{equation}
where $\alpha$, $\beta$ are Dirac indices and $i$, $j$ are indices of the
symmetry group: $i, j = 1, \ldots, d(r)$.  In analogy with electrodynamics,
we would guess that the propagator of the vector fields is
\begin{equation}
  \langle A^a_\mu(x)A^b_\nu(y)\rangle
  = \int \frac{d^4p}{(2\pi)^4}\ \frac{-ig_{\mu\nu}}{p^2+i\epsilon}\
    e^{-ip\cdot(x-y)}\,\delta^{ab} .
  \label{eq:16.5}
\end{equation}
with $a, b = 1, \ldots, d(G)$.  We will derive this formula in the next
section.

To find the vertices, we write out the nonlinear terms in \eqref{eq:16.1}.
If $\mathcal{L}_0$ is the free field Lagrangian, then
\begin{equation}
  \mathcal{L} = \mathcal{L}_0 + g\bar{\psi}\gamma^\mu t^a\psi\,A^a_\mu
    - gf^{abc}(\partial_\mu A^a_\nu)A^{b\mu}A^{c\nu}
    - \frac{1}{4}g^2\bigl( f^{eab}A^a_\mu A^b_\nu \bigr)\bigl(
      f^{ecd}A^{c\mu}A^{d\nu} \bigr) .
  \label{eq:16.6}
\end{equation}
The first of the three nonlinear terms gives the fermion-gauge boson vertex
\begin{equation}
  ig\gamma^\mu t^a ;
  \label{eq:16.7}
\end{equation}
this is a matrix that acts on the Dirac and gauge indices of the fermions.
The second nonlinear term leads to a three gauge boson vertex.  To work out
this vertex, we first choose a definite convention for the external momenta
and Lorentz and gauge indices.  A suitable convention is shown in
Fig.~\ref{fig:16.1}, with all momenta pointing inward.  Consider first
contracting the external gauge particle with momentum $k$ to the first factor
of $A^a_\nu$, the gauge particle with momentum $p$ to the second, and the
gauge particle of momentum $q$ to the third.  The derivative contributes a
factor $(-ik_\mu)$ if the momentum points into the diagram.  Then this
contribution is
\begin{equation}
  -igf^{abc}(-ik_\mu)g_{\nu\rho} .
  \label{eq:16.8}
\end{equation}
In all, there are $3!$ possible contractions, which alternate in sign
according to the total antisymmetry of $f^{abc}$.  The sum of these is
exhibited in Fig.~\ref{fig:16.1}.  Finally, the last term of \eqref{eq:16.6}
leads to a four gauge boson vertex.  Following the conventions of
Fig.~\ref{fig:16.1}, one possible contraction gives the contribution
\begin{equation}
  -ig^2f^{eab}f^{ecd}g^{\mu\nu}g^{\rho\sigma} .
  \label{eq:16.9}
\end{equation}
There are $4!$ possible contractions, of which sets of 4 are equal to one
another.  The sum of these contributions is shown in Fig.~\ref{fig:16.1}.

\begin{figure}[t]
  \centering
  \includegraphics[width=0.65\textwidth]{fig16.1.png}
  \caption{Feynman rules for fermion and gauge boson vertices of a non-Abelian
    gauge theory.}
  \label{fig:16.1}
\end{figure}

Notice that all of these vertices involve the same coupling constant $g$.
We derived the vertices, and thus the equality of the coupling constants, as
a part of our construction of the Lagrangian from the principle of
non-Abelian gauge invariance.  However, it is also possible to see the need
for this equality a posteriori, from the properties of Feynman amplitudes.

\subsection{Equality of Coupling Constants}

One property that we expect from Feynman amplitudes in non-Abelian gauge
theories is that they should satisfy Ward identities similar to those of QED.
These Ward identities express the conservation of the symmetry currents,
which follows already from the global symmetry of the theory.  In QED, the
simplest form of the Ward identity was obtained by putting external electrons
and positrons on shell.  In non-Abelian gauge theories, the gauge bosons also
carry charge and so these must also be put on shell to remove contact terms.
With all external particles on shell, the amplitude for production of a gauge
boson should obey
\begin{equation}
  \epsilon^{*\nu}(k)\ \mathcal{M}^\nu(k) = 0
  \qquad\text{for } k^\mu\epsilon_\mu = 0 ,
  \label{eq:16.10}
\end{equation}
where the polarization vector of the produced boson satisfies the on-shell
condition.  This identity is not only an indication of the local gauge
symmetry, but is physically important in its own right.  Like the photon, the
non-Abelian gauge boson has only two physical polarization states.  In QED,
the on-shell Ward identity expressed the fact that the orthogonal, unphysical
polarization states are not produced in scattering processes.  The on-shell
Ward identity will play a similar role in the non-Abelian case.

Let us check the Ward identity in a simple case, the lowest-order diagrams
contributing to fermion-antifermion annihilation into a pair of gauge bosons.
In order $g^2$, there are three diagrams, shown in Fig.~\ref{fig:16.2}.  The
first two diagrams are similar to the QED diagrams that we studied in
Section~5.5; they sum to
\begin{equation}
  \mathcal{M} = (ig)^2\ \bar{v}(p_+)\ \Bigl[ \gamma^\mu t^a\ \frac{i}{
    \slashed{p}_+ - \slashed{k}_1 - m}\ \gamma^\nu t^b
    + \gamma^\nu t^b\ \frac{i}{\slashed{p}_+ - \slashed{k}_2 - m}\
      \gamma^\mu t^a \Bigr]\ u(p_-)\ \epsilon^{*}_\mu(k_1)
        \epsilon^{*}_\nu(k_2) .
  \label{eq:16.11}
\end{equation}
The vectors $\epsilon(k_i)$ are the gauge boson polarization vectors; for
physical polarizations, these satisfy $k^\mu_i\epsilon_\mu(k_i) = 0$.  To check
the Ward identity \eqref{eq:16.10}, we replace $\epsilon^{*}_\nu(k_2)$ in
\eqref{eq:16.11} by $k_{2\nu}$.  This gives
\begin{equation}
  \begin{aligned}
    \mathcal{M} &= (ig)^2\ \bar{v}(p_+)\ \Bigl[ \gamma^\mu t^a\ \frac{\slashed{k}_2}
      {\slashed{p}_+ - \slashed{k}_1 - m}\ \gamma^\nu t^b
      + \gamma^\nu t^b\ \frac{\slashed{k}_2}{\slashed{p}_+ - \slashed{k}_2 - m}\
        \gamma^\mu t^a \Bigr]\ u(p_-)\ \epsilon^{*}_\mu(k_1) .
  \end{aligned}
  \label{eq:16.12}
\end{equation}
Since
\begin{equation}
  (\slashed{p}_- - m)u(p_-) = 0
  \qquad\text{and}\qquad
  \bar{v}(p_+)(\slashed{p}_+ - m) = 0 ,
  \label{eq:16.13}
\end{equation}
we can add these quantities to $\slashed{k}_2$ in the first and second terms
of \eqref{eq:16.12}, to cancel the denominators.  This gives
\begin{equation}
  \mathcal{M} = (ig)^2\ \bar{v}(p_+)\ \Bigl[ -\frac{i}{-i}\ \gamma^\mu\,
    \frac{1}{2}\bigl[ t^a, t^b \bigr] \Bigr]\ u(p_-)\ \epsilon^{*}_\mu(k_1) .
  \label{eq:16.14}
\end{equation}
In the Abelian case, this expression would vanish.  In the non-Abelian case,
however, the residual term is nonzero and depends on the commutator of gauge
group generators:
\begin{equation}
  \mathcal{M} = -\frac{g^2}{2}\ \bar{v}(p_+)\gamma^\mu u(p_-)\ \epsilon^{*}_\mu(k_1)
    \cdot f^{abc}t^c .
  \label{eq:16.15}
\end{equation}
We need to find another contribution to cancel this term.  Notice, however,
that this term has the group index structure of a fermion-gauge boson vertex
$(g\bar{\psi}\gamma^\mu t^c\psi)$ multiplied by a three gauge boson vertex
$(gf^{abc})$.  This is just the structure of the third diagram in
Fig.~\ref{fig:16.2}.

\begin{figure}[t]
  \centering
  \includegraphics[width=0.6\textwidth]{fig16.2.png}
  \caption{Diagrams contributing to fermion-antifermion annihilation to two
    gauge bosons.}
  \label{fig:16.2}
\end{figure}

To check that the cancellation works, let us evaluate the third diagram:
\begin{equation}
  \mathcal{M} = ig\ \bar{v}(p_+)\gamma^\rho t^c u(p_-)\ \frac{1}{k_3^2}\
    \epsilon^{*}_\mu(k_1)\epsilon^{*}_\nu(k_2)\ gf^{abc}\ \bigl[
      g^{\mu\nu}(k_2-k_1)^\rho + g^{\nu\rho}(k_3-k_2)^\mu
        + g^{\rho\mu}(k_1-k_3)^\nu \bigr] ,
  \label{eq:16.16}
\end{equation}
with $k_3 = -k_1 - k_2$.  If we replace $\epsilon^{*}_\nu(k_2)$ with
$k_{2\nu}$, then eliminate $k_2$ using momentum conservation, the expression
in brackets simplifies as follows:
\begin{equation}
  \epsilon^{*}_\mu(k_1)\ \bigl[ g^{\mu\nu}(k_2-k_1)^\rho + g^{\nu\rho}
    (k_3-k_2)^\mu + g^{\rho\mu}(k_1-k_3)^\nu \bigr]
  \to \epsilon^{*}_\mu(k_1)\ \bigl[ k_2^\mu(k_2-k_1)^\rho + k_2^\rho
    (k_3-k_2)^\mu + g^{\rho\mu}(k_1-k_3)\cdot k_2 \bigr] ,
  \label{eq:16.17}
\end{equation}
which combines with \eqref{eq:16.15} to give zero.  Thus the Ward identity
is satisfied only if the coupling constant in the three-gauge-boson vertex is
equal to that in the fermion vertex.  This argument can be generalized to
show that all of the vertices of a non-Abelian gauge theory share the same
coupling constant $g$.

\section{The Faddeev-Popov Lagrangian}

It is not surprising that we have found a problem with our Feynman rules for
non-Abelian gauge theories, since we were not very careful in deriving them.
In particular, we did not actually derive expression \eqref{eq:16.5} for the
gauge field propagator.  In this section we will remedy this by going through
a formal derivation of this expression.  We will find that, although
expression \eqref{eq:16.5} is indeed correct, it is incomplete: It must be
supplemented by additional rules of a completely new type.

To define the functional integral for a theory with non-Abelian gauge
invariance, we will use the Faddeev-Popov method, as introduced in
Section~9.4 to quantize the electromagnetic field.  Our present discussion
will follow Section~9.4 closely.  However, as we have by now come to expect,
the case of non-Abelian local symmetry brings with it new tricks and
surprises.

First consider the quantization of the pure gauge theory, without fermions.
To derive the Feynman rules, we must define the functional integral
\begin{equation}
  Z = \int \mathcal{D}A\ e^{iS[A]} .
  \label{eq:16.22}
\end{equation}
As in the Abelian case, the Lagrangian is unchanged along the infinite number
of directions in the space of field configurations corresponding to local
gauge transformations.  To compute the functional integral we must factor out
the integrations along these directions, constraining the remaining integral
to a much smaller space.

As in electrodynamics, we will constrain the gauge directions by applying a
gauge-fixing condition $G^a(A) = 0$ at each point $x$.  Following Faddeev
and Popov, we can introduce this constraint by inserting into the functional
integral the identity \eqref{eq:9.53}:
\begin{equation}
  1 = \Delta[A]\ \int \mathcal{D}\alpha\ \prod_{x,a} \delta\bigl( G^a(A^\alpha)
    \bigr) ,
  \label{eq:16.23}
\end{equation}
where $A^\alpha$ is the gauge field $A$ transformed through a finite gauge
transformation as in \eqref{eq:15.47}:
\begin{equation}
  A^{a\alpha}_\mu = A^a_\mu + \frac{1}{g}\partial_\mu\alpha^a
    + f^{abc}A^b_\mu\alpha^c .
  \label{eq:16.24}
\end{equation}
In evaluating the determinant, the infinitesimal form of this transformation
will be more useful:
\begin{equation}
  (A^\alpha)^a_\mu = A^a_\mu + \partial_\mu\alpha^a + f^{abc}A^b_\mu\alpha^c
  = A^a_\mu + D_\mu\alpha^a ,
  \label{eq:16.25}
\end{equation}
where $D_\mu$ is the covariant derivative \eqref{eq:15.86} acting on a field
in the adjoint representation.  Note that, as long as the gauge-fixing
function $G^a(A)$ is linear, the functional derivative
$\delta G^a(A^\alpha)/\delta\alpha$ is independent of $\alpha$.

Since the Lagrangian is gauge invariant, we can replace $A$ by $A^\alpha$ in
the exponential of \eqref{eq:16.22}.  Then, as in the Abelian case, we can
interchange the order of the functional integrals over $A$ and $\alpha$, and
then change variables in the inner integral from $A$ to $A' = A^\alpha$.  The
transformation \eqref{eq:16.24} looks more complicated than in the Abelian
case, but it is nothing more than a linear shift of the $A^a_\mu$, followed by
a unitary rotation of the various components of the symmetry multiplet
$A^a_\mu(x)$ at each point.  Both of these operations preserve the measure
\begin{equation}
  \mathcal{D}A = \prod_{x,a,\mu} dA^a_\mu .
  \label{eq:16.26}
\end{equation}
Thus $\mathcal{D}A = \mathcal{D}A'$, under the integral over $\alpha$.  Just
as in the Abelian case, the integral over gauge motions $\alpha$ can be
factored out of the functional integral into an overall normalization, leaving
us with
\begin{equation}
  \int \mathcal{D}A\ e^{iS[A]} = \Bigl( \int \mathcal{D}\alpha \Bigr)\
    \int \mathcal{D}A\ e^{iS[A]}\ \delta\bigl( G(A) \bigr)\ \det\Bigl(
      \frac{\delta G}{\delta\alpha} \Bigr) .
  \label{eq:16.27}
\end{equation}
This normalization factor cancels in the computation of correlation functions
of gauge-invariant operators.

From this point, the derivation of the gauge boson propagator proceeds as for
the photon propagator.  We choose the generalized Lorentz gauge condition
\begin{equation}
  G^a(A) = \partial^\mu A^a_\mu(x) - \omega^a(x) ,
  \label{eq:16.28}
\end{equation}
with a Gaussian weight for $\omega^a$ as in Eq.~\eqref{eq:9.56}.  The
manipulations of Section~9.4 then lead to the class of gauge field
propagators
\begin{equation}
  \langle A^a_\mu(x)A^b_\nu(y)\rangle = \int \frac{d^4p}{(2\pi)^4}\
    \frac{-i}{p^2+i\epsilon}\ \Bigl[ g_{\mu\nu} - (1-\xi)\frac{p_\mu p_\nu}
      {p^2} \Bigr]\ e^{-ip\cdot(x-y)}\,\delta^{ab} ,
  \label{eq:16.29}
\end{equation}
with a freely adjustable gauge parameter $\xi$.  Our guess \eqref{eq:16.5}
corresponds to the choice $\xi = 1$, called the \emph{Feynman-'t Hooft
gauge}.

So far, this whole derivation parallels the case of electrodynamics.  Here,
however, there is one more nontrivial ingredient.  In QED, the determinant in
Eq.~\eqref{eq:16.23} was independent of $A$, so this quantity could be
treated as just another contribution to the normalization factor.  In the
non-Abelian case this is no longer true.  Using the infinitesimal form
\eqref{eq:16.25} of the gauge transformation, we can evaluate
\begin{equation}
  \frac{\delta G^a(A^\alpha)}{\delta\alpha^b} = \frac{1}{g}\partial^2\delta^{ab}
    - \frac{1}{g}\partial^\mu f^{abc}A^c_\mu ,
  \label{eq:16.30}
\end{equation}
acting on a field in the adjoint representation; this operator depends on
$A$.  The functional determinant of \eqref{eq:16.30} thus contributes new
terms to the Lagrangian.

Faddeev and Popov chose to represent this determinant as a functional
integral over a new set of anticommuting fields belonging to the adjoint
representation:
\begin{equation}
  \det\Bigl( -\frac{1}{g}\partial^\mu D_\mu \Bigr)
  = \int \mathcal{D}c\,\mathcal{D}\bar{c}\ \exp\Bigl[ i\int d^4x\
    \bar{c}^a\Bigl( -\frac{1}{g}\partial^\mu D_\mu \Bigr)^{ab}c^b \Bigr] .
  \label{eq:16.31}
\end{equation}
We derived this formal identity in Eq.~(9.69), using our rules for
fermionic functional integrals.  (The factor of $1/g$ is absorbed into the
normalization of the fields $c$ and $\bar{c}$.)  But to give the correct
identity, $c$ and $\bar{c}$ must be anticommuting fields that are scalars
under Lorentz transformations.  The quantum excitations of these fields have
the wrong relation between spin and statistics to be physical particles.
However, we can nevertheless treat these excitations as additional particles
in the computation of Feynman diagrams.  These new fields and their particle
excitations are called \emph{Faddeev-Popov ghosts}.

If we temporarily suppress our curiosity about the physical interpretation of
the ghosts, we can work out their Feynman rules.  We write the ghost
Lagrangian more explicitly as
\begin{equation}
  \mathcal{L}_{\mathrm{ghost}} = \bar{c}^a\bigl( -\partial^2\delta^{ac}
    - g\partial^\mu f^{abc}A^b_\mu \bigr)c^c .
  \label{eq:16.32}
\end{equation}
The first term gives a ghost propagator,
\begin{equation}
  \langle c^a(x)\bar{c}^b(y)\rangle = \int \frac{d^4p}{(2\pi)^4}\
    \frac{i\delta^{ab}}{p^2+i\epsilon}\ e^{-ip\cdot(x-y)} .
  \label{eq:16.33}
\end{equation}
In a diagram, this propagator carries an arrow that shows the flow of ghost
number, as in Fig.~\ref{fig:16.5}.  In the interaction term of \eqref{eq:16.32},
the derivative stands to the left of the gauge field; this implies that this
derivative is evaluated with the momentum coming out of the vertex along the
ghost line.  The explicit Feynman rule is shown in Fig.~\ref{fig:16.5}.  As
with the other vertices we have encountered, the coupling constant $g$ that
appears in this vertex must be equal to the coupling constant $g$ in the
three-boson vertex in order to avoid upsetting the Ward identities.

\begin{figure}[t]
  \centering
  \includegraphics[width=0.5\textwidth]{fig16.5.png}
  \caption{Feynman rules for Faddeev-Popov ghosts.}
  \label{fig:16.5}
\end{figure}

There are no further subtleties in the construction of the perturbation
theory for non-Abelian gauge theories.  In particular, it is straightforward
to include fermions.  The final Lagrangian, including all of the effects of
Faddeev-Popov gauge fixing, is
\begin{equation}
  \mathcal{L} = -\frac{1}{4}(F^a_{\mu\nu})^2 - \frac{1}{2\xi}
    (\partial^\mu A^a_\mu)^2 + \bar{\psi}(i\slashed{D} - m)\psi
    + \bar{c}^a(-\partial^\mu D_\mu)^{ab}c^b .
  \label{eq:16.34}
\end{equation}
This Lagrangian leads to the propagator \eqref{eq:16.29}, and to the set of
Feynman rules for vertices shown in Figs.~\ref{fig:16.1} and
\ref{fig:16.5}.

The argument we have just completed suffices to derive the Feynman diagram
expansion of any correlation function of gauge-invariant operators in a
non-Abelian gauge theory.  At the end of Section~9.4, we explained that the
Faddeev-Popov gauge-fixing technique also gives the correct gauge-invariant
expressions for $S$-matrix elements.  This remains true in the non-Abelian
case.  However, the argument given in Section~9.4 relied upon the
cancellation in QED of the emission probabilities for timelike and
longitudinal photons, and we have already found that this cancellation does
not go through in the non-Abelian case.  In Section~16.4 we will construct a
more sophisticated argument, in which the Faddeev-Popov ghosts play an
essential role, that will correctly generalize our previous argument to
non-Abelian gauge theories.

\section{Ghosts and Unitarity}

We might now ask whether the new ingredients that we found in the previous
section, the Faddeev-Popov ghosts, can resolve the paradox that we
encountered at the end of Section~16.1.  There we saw that the first diagram
in Fig.~\ref{fig:16.6} contains a nonzero contribution to its imaginary part
that does not correspond to a possible final state with physical gauge boson
polarizations.  We will now compute this contribution more carefully.  We
must then add a new potential contribution from the ghosts, shown as the
second diagram in Fig.~\ref{fig:16.6}.

\begin{figure}[t]
  \centering
  \includegraphics[width=0.55\textwidth]{fig16.6.png}
  \caption{The diagram on the left, in which each circle represents the sum
    of the three contributions of Fig.~\ref{fig:16.2}, gives a possible
    problem for the optical theorem.  The ghost diagram on the right cancels
    the anomalous terms.}
  \label{fig:16.6}
\end{figure}

Let us call the amplitude for fermion-fermion annihilation into gauge bosons,
which we studied in Section~16.1,
\begin{equation}
  \mathcal{M}(p_+, p_-; k_1, k_2) ,
  \label{eq:16.35}
\end{equation}
the amplitude for two gauge bosons to convert to a fermion-antifermion pair
will be, correspondingly, $\mathcal{M}'$.  Then, following the Cutkosky rules
of Section~7.3, we find the imaginary part of the first diagram in
Fig.~\ref{fig:16.6} by replacing the cut gauge boson propagator with momentum
$k_i$ by
\begin{equation}
  (-ig_{\mu\nu}\delta^{ab})\cdot(-2\pi i)\delta(k_i^2) .
  \label{eq:16.36}
\end{equation}
Replacing both propagators gives two delta functions, turning the
four-dimensional integrals over the gauge boson momenta into three-dimensional
phase space integrals, as in the example in Section~7.3.  We are thus left
with the expression
\begin{equation}
  \frac{1}{2}\ \int d\Pi_2\ \mathcal{M}\ \mathcal{M}' ,
  \label{eq:16.37}
\end{equation}
integrated over the phase space of two massless particles.  The factor $1/2$
is a symmetry factor for the Feynman diagram or, equivalently, a correction
to the phase space integral for identical particles.

Now introduce the representation Eq.~(16.20) for $g_{\mu\rho}$ and
$g_{\nu\sigma}$.  The pieces that involve only transverse polarizations
correspond to the expected imaginary parts necessary to satisfy the optical
theorem.  We need not consider these terms further.  The cross terms between
physical and unphysical polarizations vanish: We showed in Section~16.1 that
\begin{equation}
  i\mathcal{M}^\mu\,\epsilon^{*}_{\mu}(k_1)\,\epsilon^{*}_{\nu}(k_2)
    \cdot\cdots = 0 .
  \label{eq:16.38}
\end{equation}
The same identity holds if $\mathcal{M}$ is replaced by $\mathcal{M}'$, and
if $\epsilon^{+}$ is replaced by $\epsilon^{-}$.  Furthermore, the amplitude
vanishes if both polarization vectors are forward or both are backward.  The
only surviving terms are the cross terms between forward and backward
polarization, which yield the expression
\begin{equation}
  \int d\Pi_2\ \mathcal{M}^\mu\epsilon^{+}_{\mu}\ \mathcal{M}'^\nu
    \epsilon^{-}_{\nu}\ \cdots ,
  \label{eq:16.39}
\end{equation}
integrated over phase space.  We worked out the value of the first factor in
Eq.~(16.21), and the contraction with $\mathcal{M}'$ is very
similar.  Substituting these results, expression \eqref{eq:16.39} becomes
\begin{equation}
  \int d\Pi_2\ (ig\bar{v}(p_+)\gamma^\rho t^c u(p_-))\cdot g f^{abd}\ \cdots .
  \label{eq:16.40}
\end{equation}
Using the identity
\begin{equation}
  \bar{v}(p_+)\slashed{k}_1(\slashed{k}_1+\slashed{k}_2)u(p_-)
  = \bar{v}(p_+)\slashed{k}_1(\slashed{p}_+ + \slashed{p}_-)u(p_-) = 0 ,
  \label{eq:16.41}
\end{equation}
we see that the two terms added in \eqref{eq:16.40} are equal.

Now add the contribution from the Faddeev-Popov ghosts.  Using the Feynman
rules in Fig.~\ref{fig:16.5}, we can assemble the amplitude for
fermion-antifermion annihilation into a pair of ghosts:
\begin{equation}
  i\mathcal{M}_{\mathrm{ghost}} = ig\bar{v}(p_+)\gamma^\mu t^c u(p_-)\
    \frac{1}{(k_1+k_2)^2}\ \cdot\ gf^{abc}k^5_2\ \cdots .
  \label{eq:16.42}
\end{equation}
Combining the ghost contribution with \eqref{eq:16.40}, one finds that the
unphysical contributions cancel exactly, restoring the optical theorem.  This
is a first indication that the Faddeev-Popov ghosts play an essential role in
the unitarity of non-Abelian gauge theories.  In the next section we will
develop a more general formalism, based on the BRST symmetry, that makes this
cancellation manifest to all orders.

\section{BRST Symmetry}

The gauge-fixing procedure of Section~16.2 introduced new fields into the
theory: the ghosts $c^a$, the antighosts $\bar{c}^a$, and the auxiliary
fields $B^a$ (in the Faddeev-Popov construction, the gauge-fixing function
introduces a field $B^a$ that is the Lagrange multiplier enforcing the gauge
condition).  The full gauge-fixed Lagrangian is invariant under a new, global
symmetry, called the \emph{BRST symmetry}, after Becchi, Rouet, Stora, and
Tyutin.  This symmetry relates the ghosts, the gauge fields, and the
fermions in a nontrivial way.  It is generated by a nilpotent charge $Q$,
which plays a central role in the proof of unitarity.

The BRST transformation acts on the fields as follows:
\begin{equation}
  \begin{aligned}
    \delta A^a_\mu &= \epsilon D_\mu c^a;\\
    \delta\psi &= ig\epsilon\,c^a t^a\psi;\\
    \delta c^a &= -\frac{1}{2}g\epsilon f^{abc}c^b c^c;\\
    \delta\bar{c}^a &= -i\epsilon B^a;\\
    \delta B^a &= 0 .
  \end{aligned}
  \label{eq:16.45}
\end{equation}
Here $\epsilon$ is a Grassmann-valued transformation parameter.  The
transformation is nilpotent in the sense that applying it twice gives zero:
\begin{equation}
  \delta(\delta\Phi) = 0
  \label{eq:16.48}
\end{equation}
for every field $\Phi$.  Corresponding to this continuous symmetry, there is
a conserved charge $Q$ that commutes with the Hamiltonian.  The relation
\eqref{eq:16.48} is equivalent to the operator identity
\begin{equation}
  Q^2 = 0 .
  \label{eq:16.50}
\end{equation}
We say that the BRST operator $Q$ is \emph{nilpotent}.

A nilpotent operator that commutes with $H$ divides the eigenstates of $H$
into three subspaces.  Many eigenstates of $H$ must be annihilated by $Q$ so
that \eqref{eq:16.50} can be satisfied.  Let $\mathcal{H}_1$ be the subspace
of states that are not annihilated by $Q$.  Let $\mathcal{H}_2$ be the
subspace of states of the form
\begin{equation}
  |\psi_2\rangle = Q|\psi_1\rangle ,
  \label{eq:16.51}
\end{equation}
where $|\psi_1\rangle$ is in $\mathcal{H}_1$.  According to \eqref{eq:16.50},
acting $Q$ again on these states gives zero.  Finally, let $\mathcal{H}_0$ be
the subspace of states $|\psi_0\rangle$ that satisfy $Q|\psi_0\rangle = 0$
but that cannot be written in the form \eqref{eq:16.51}.  The subspace
$\mathcal{H}_2$ is quite peculiar, because any two states in this subspace
have zero inner product:
\begin{equation}
  \langle\psi_{2a}|\psi_{2b}\rangle = \langle\psi_{1a}|Q|\psi_{2b}\rangle
  = 0
  \label{eq:16.52}
\end{equation}
by \eqref{eq:16.50}.  By the same argument, the states of $\mathcal{H}_2$
have zero inner product with the states of $\mathcal{H}_0$.

These considerations seem extremely abstract, but they have a direct physical
correspondence.%
\footnote{The following argument is presented only at an intuitive level.
For a rigorous discussion, see T.~Kugo and I.~Ojima, \emph{Suppl.\ Prog.\
Theor.\ Phys.} \textbf{66}, 1 (1979).}  To see this, consider single-particle
states of the non-Abelian gauge theory in the limit of zero coupling.
According to the transformation \eqref{eq:16.45}, $Q$ converts the forward
component of $A^a_\mu$ to a ghost field; equivalently, $Q$ converts a single
forward-polarized gauge boson to a ghost.  At $g = 0$, $Q$ annihilates the
one-ghost state.  At the same time, $Q$ converts the antighost state to a
quantum of $B^a$.  To identify this state, note that the Lagrangian
\eqref{eq:16.34} implies the classical field equation
\begin{equation}
  \frac{1}{\xi}\partial^\mu A^a_\mu = -B^a .
  \label{eq:16.53}
\end{equation}
Thus the quanta of the field $B^a$ are those quanta of $A^a_\mu$ with
polarization vectors such that $k^\mu\epsilon_\mu(k) \neq 0$; these are the
backward-polarized gauge bosons.

We have now seen that, among the single-particle states of the gauge theory,
forward gauge bosons and antighosts belong to $\mathcal{H}_1$, ghosts and
backward gauge bosons belong to $\mathcal{H}_2$, and transverse gauge bosons
belong to $\mathcal{H}_0$.  More generally, it can be shown that asymptotic
states containing ghosts, antighosts, or gauge bosons of unphysical
polarization always belong to $\mathcal{H}_1$ or $\mathcal{H}_2$, while the
asymptotic states in $\mathcal{H}_0$ are those with only transversely
polarized gauge bosons.  The BRST operator thus gives a precise relation
between the unphysical gauge boson polarization states and the ghosts and
antighosts as positive and negative degrees of freedom.

In Section~9.4, we argued that the Faddeev-Popov prescription gave the
correct, gauge-invariant result for a certain subclass of $S$-matrix
elements, from which we could compute the physical scattering cross sections
of transversely polarized gauge bosons.  These $S$-matrix elements were
constructed by putting operators in the far past to create transversely
polarized gauge bosons, adiabatically turning on the gauge coupling,
adiabatically turning off the gauge coupling, and then placing operators in
the far future to annihilate gauge bosons with transverse polarization.
However, this argument had a possible problem: If the states created as
collections of transversely polarized bosons in the far past could evolve
into states that contained gauge bosons of other polarizations in the far
future, the $S$-matrix projected between transverse gauge boson states would
not be unitary.  This problem would also lead to the technical problem
discussed in the previous section: The Cutkosky cuts of diagrams contributing
to $S$-matrix elements would have nonzero contributions from unphysical
polarizations.  In Section~9.4, we used an argument special to the Abelian
case to show that these problems do not arise in QED.  In the non-Abelian
case, the removal of unphysical gauge boson polarizations is more subtle, and
we have seen that it involves the ghosts in an essential way.

To resolve this subtle problem, we apply the principle of BRST symmetry.
Let $|A;\mathrm{tr}\rangle$ be an external state that contains no ghosts or
antighosts and only gauge bosons with transverse polarization.  We wish to
show that the $S$-matrix projected onto such states is unitary:
\begin{equation}
  \sum_C \langle A;\mathrm{tr}|S^{\dagger}|C;\mathrm{tr}\rangle\,
    \langle C;\mathrm{tr}|S|B;\mathrm{tr}\rangle
  = \langle A;\mathrm{tr}|1|B;\mathrm{tr}\rangle .
  \label{eq:16.54}
\end{equation}
As we explained above, the physical states $|A;\mathrm{tr}\rangle$ belong
to---and, in fact, span---the subspace $\mathcal{H}_0$ defined by the BRST
operator.  In particular, all of these states are annihilated by $Q$.  Since
$Q$ commutes with the Hamiltonian, the time evolution of any such state must
also produce a state annihilated by $Q$.  Thus,
\begin{equation}
  Q\cdot S|A;\mathrm{tr}\rangle = 0 .
  \label{eq:16.55}
\end{equation}
This implies that the states $S|A;\mathrm{tr}\rangle$ must be linear
combinations of states in $\mathcal{H}_0$ and $\mathcal{H}_2$.  However,
states in $\mathcal{H}_2$ have zero inner product with one another and with
states in $\mathcal{H}_0$.  Thus the inner product of any two states of the
form $S|A;\mathrm{tr}\rangle$ comes only from the overlap of the components
in $\mathcal{H}_0$, so we can write
\begin{equation}
  \langle A;\mathrm{tr}|S^{\dagger}\cdot S|B;\mathrm{tr}\rangle
  = \sum_C \langle A;\mathrm{tr}|S^{\dagger}|C;\mathrm{tr}\rangle\,
    \langle C;\mathrm{tr}|S|B;\mathrm{tr}\rangle .
  \label{eq:16.56}
\end{equation}
Since the full $S$-matrix is unitary, this relation implies that the
restricted $S$-matrix is also unitary, Eq.~\eqref{eq:16.54}.  In addition,
\eqref{eq:16.56} implies that the sum of the Cutkosky cuts of diagrams
contributing to the $S$-matrix in a given order is equal to the sum of the
cuts involving transverse gauge bosons only.  Thus, the cancellation between
diagrams that produce pairs of gauge bosons with unphysical polarizations and
those that produce ghosts is a general property that persists to all orders
in perturbation theory.

Since the BRST transformation generates a continuous symmetry, it generates
a set of Ward identities.  These identities are similar in structure to the
Ward identities of the non-Abelian gauge symmetry, since the BRST symmetry
contains a gauge transformation whose parameter is the ghost field.  However,
the identities that follow from BRST symmetry are simpler.  We will not study
the Ward identities of non-Abelian gauge theory further in this book.
However, when one discusses the renormalization of gauge theories at a higher
level, the central identities among renormalization constants that follow
from the Ward identities are most easily derived using the BRST symmetry.%
\footnote{An introduction to the Ward identities of the BRST symmetry is
given by Taylor (1976).}

\section{One-Loop Divergences of Non-Abelian Gauge Theory}

Now that we have discussed the general properties of tree-level diagrams in
non-Abelian gauge theories, we turn our attention to diagrams with loops.  As
always in quantum field theory, some of these loop diagrams will diverge, and
we must take care to treat the divergent integrals correctly.

The Lagrangian of a non-Abelian gauge theory \eqref{eq:15.39} contains no
interactions of dimension higher than 4.  Therefore, by the general arguments
of Chapter~10, this Lagrangian is renormalizable, in the sense that the
divergences can be removed by a finite number of counterterms.  However, in
non-Abelian gauge theories, as in QED, the gauge symmetries of the theory
imply stronger restrictions on the structure of the divergences.  In QED,
provided that we use a gauge-invariant regulator, there are only four possible
divergent coefficients, which are subtracted by the counterterms for the
electromagnetic vertex ($\delta_1$), for the electron and photon field
strength ($\delta_2$ and $\delta_3$), and for the electron mass ($\delta_m$).
In particular, the possibility of a photon mass renormalization is excluded by
gauge invariance.  Furthermore, the two counterterms $\delta_1$ and
$\delta_2$ are equal to one another, and cancel in the evaluation of the
electron-photon vertex function, as a consequence of the Ward identity.
Non-Abelian gauge symmetries imply similar restrictions on the divergences of
Feynman diagrams.  In this section, we will illustrate some of these
restrictions through examples of one-loop diagrams.

\subsection{The Gauge Boson Self-Energy}

In QED, the strongest constraints of gauge invariance come in the evaluation
of the photon self-energy.  The Ward identity implies the relation
\begin{equation}
  q_\mu\Pi^{\mu\nu}(q) = 0 ,
  \label{eq:16.57}
\end{equation}
which in turn implies that the photon self-energy diagrams have the structure
\begin{equation}
  \Pi^{\mu\nu}(q) = (g^{\mu\nu}q^2 - q^\mu q^\nu)\,\Pi(q^2) .
  \label{eq:16.58}
\end{equation}
The only divergence possible is a logarithmically divergent contribution to
$\Pi(q^2)$.  In non-Abelian gauge theories, \eqref{eq:16.57} still holds, so
the self-energy again has the Lorentz structure \eqref{eq:16.58}.  However,
the cancellations that lead to this structure are more complex.  Here we will
exhibit these cancellations by computing the gauge boson self-energy in
detail at the one-loop level.  In order to preserve gauge invariance, we will
use dimensional regularization.

\begin{figure}[t]
  \centering
  \includegraphics[width=0.6\textwidth]{fig16.7.png}
  \caption{Contributions to the gauge boson self-energy in order $g^2$.}
  \label{fig:16.7}
\end{figure}

The contributions of order $g^2$ to the gauge boson self-energy are shown in
Fig.~\ref{fig:16.7}.  (In addition to these 1PI diagrams, there are three
``tadpole'' diagrams; but these automatically vanish, as in QED, by the
argument given below Eq.~\eqref{eq:10.5}.)  The fermion loop diagram can be
considered separately from the other diagrams, since in principle we could
include any number of fermions in the theory.  We will see below that the
contributions of the three remaining diagrams interlock in an essential way.

Let us first calculate the fermion loop diagram.  The Feynman rule for the
vertices in this diagram is identical to the QED Feynman rule, except for the
addition of a group matrix $t^a$ that acts on the fermion gauge group indices.
The value of this diagram is therefore the same as in QED, Eq.~(7.90),
multiplied by a trace over group matrices:
\begin{equation}
  \text{(fermion loop)} \propto \mathrm{tr}[t^at^b]\,\Pi^{\mu\nu}_{\mathrm{QED}} .
\end{equation}
The value of the trace is given by Eq.~\eqref{eq:15.78}:
$\mathrm{tr}[t^at^b] = C(r)\delta^{ab}$.  In a theory with several species of
fermions, there would be a diagram of this type for each species.  We will be
mainly interested in the divergent part of this diagram, which is independent
of the fermion mass.  If there are $n_f$ species of fermions, all in the same
representation $r$, then the total contribution of fermion loop diagrams takes
the form
\begin{equation}
  \Bigl( \sum_{\text{fermions}} \Bigr)^{\mu\nu}(q)
  = i(g^{\mu\nu}q^2 - q^\mu q^\nu)\ \frac{g^2}{(4\pi)^2}\ n_f\,C(r)\,
    \Gamma(2-d/2)\ +\ \cdots .
  \label{eq:16.59}
\end{equation}

Now consider the three diagrams from the pure gauge sector.  The contribution
of these diagrams depends on the gauge; we will use Feynman-'t Hooft gauge,
$\xi = 1$.  Using the three-gauge-boson vertex from Fig.~\ref{fig:16.1}, we
can write the first of the three diagrams as
\begin{equation}
  \begin{aligned}
    \Pi^{\mu\nu}_{(1)}(q) &= -\frac{1}{2}\int \frac{d^dp}{(2\pi)^d}\
      \frac{1}{p^2}\frac{1}{(p+q)^2}\ g^2\ f^{acd}f^{bcd}\ N^{\mu\nu}(p,q)\\
    &= -\frac{1}{2}g^2C_2(G)\delta^{ab}\ \int \frac{d^dp}{(2\pi)^d}\
      \frac{N^{\mu\nu}(p,q)}{p^2(p+q)^2} ,
  \end{aligned}
  \label{eq:16.60}
\end{equation}
where the numerator structure is
\begin{equation}
  \begin{aligned}
    N^{\mu\nu}(p,q) &= \bigl[ g^{\mu\rho}(q-p)^\sigma + g^{\rho\sigma}(2p+q)^\mu
      + g^{\sigma\mu}(-p-2q)^\rho \bigr]\\
    &\quad \times \bigl[ g^{\nu}_{\ \rho}(p-q)_\sigma + g_{\rho\sigma}(-2p-q)^\nu
      + g^\nu_{\ \sigma}(p+2q)_\rho \bigr] .
  \end{aligned}
\end{equation}
The overall factor of $1/2$ is a symmetry factor.  The contraction of
structure constants can be evaluated using Eq.~\eqref{eq:15.93}:
$f^{acd}f^{bcd} = C_2(G)\delta^{ab}$.

To simplify the expression further, combine denominators in the standard way:
\begin{equation}
  \frac{1}{p^2(p+q)^2} = \int_0^1 dx\ \frac{1}{\bigl( (1-x)p^2 + x(p+q)^2
    \bigr)^2} = \int_0^1 dx\ \frac{1}{(P^2-\Delta)^2} ,
  \label{eq:16.61}
\end{equation}
where $P = p + xq$ and $\Delta = -x(1-x)q^2$.  Then \eqref{eq:16.60} can be
rewritten
\begin{equation}
  \Pi^{\mu\nu}_{(1)}(q) = -\frac{1}{2}g^2C_2(G)\delta^{ab}\ \int_0^1 dx\
    \int \frac{d^dP}{(2\pi)^d}\ \frac{N^{\mu\nu}(P,q)}{(P^2-\Delta)^2} .
  \label{eq:16.62}
\end{equation}
The numerator structure can be simplified by eliminating $p$ in favor of
$P$, discarding terms linear in $P$ (which integrate symmetrically to zero),
and using the standard integral identities.  Carrying out the algebra, the
sum of the three pure-gauge diagrams simplifies to a manifestly transverse
expression:
\begin{equation}
  \Pi^{\mu\nu}_{\mathrm{gauge}}(q) = (g^{\mu\nu}q^2 - q^\mu q^\nu)\,
    \Pi_{\mathrm{gauge}}(q^2) ,
  \label{eq:16.70}
\end{equation}
as required by the Ward identity of the non-Abelian gauge theory.  For future
reference, we record the ultraviolet divergent part of \eqref{eq:16.70}:
\begin{equation}
  \Pi^{\mu\nu}_{\mathrm{gauge}}(q) \supset (g^{\mu\nu}q^2 - q^\mu q^\nu)\,
    \frac{g^2}{(4\pi)^2}\ \frac{10}{3}\,C_2(G)\ \Gamma(2-d/2) .
  \label{eq:16.71}
\end{equation}

As we noted above, the result \eqref{eq:16.70} depends on the gauge used in
the calculation.  In any gauge, the boson self-energy is transverse and free
of quadratic divergences.  However, the coefficient of the transverse Lorentz
structure may depend on $\xi$.  It turns out that, for a general value of
$\xi$, the coefficient of the ultraviolet divergence in \eqref{eq:16.71} is
modified according to
\begin{equation}
  \Pi_{\mathrm{gauge}}(q^2) \sim \frac{g^2}{(4\pi)^2}\ \Bigl( \frac{10}{3}
    - \frac{\xi-1}{2} \Bigr)\,C_2(G)\ \Gamma(2-d/2) .
  \label{eq:16.72}
\end{equation}
The fact that the boson self-energy depends on the gauge does not contradict
the general theorem that $S$-matrix elements are independent of $\xi$.  The
full set of one-loop corrections to a gauge theory $S$-matrix element always
involves a number of different radiative corrections to vertices and
propagators; the gauge dependence cancels in an intricate fashion among these
various terms.

\subsection{The $\beta$ Function}

The simplest calculation that involves a gauge-invariant combination of
radiative corrections is the computation of the leading term of the
Callan-Symanzik $\beta$ function of a non-Abelian gauge theory.  The
invariance of the leading term of $\beta$ could be argued intuitively, by
saying that the coupling constant of the gauge theory should not evolve to
large values in one scheme of calculation while it stays small in another
scheme.  In Section~17.2 we will demonstrate this result more cleanly by
showing that the leading coefficient of the $\beta$ function can be extracted
from a physical cross section and so must be gauge independent.  (Surprisingly,
this conclusion actually applies to the first two coefficients of the $\beta$
function, written as a power series in $g$.)

Recall from Section~12.2 that the $\beta$ function gives the rate at which
the renormalized coupling constant changes as the renormalization scale $M$ is
increased.  Since Green's functions depend on $M$ through the counterterms
that subtract ultraviolet divergences, $\beta$ can be computed from the
counterterms that enter an appropriately chosen Green's function.  For
example, in Eq.~\eqref{eq:12.58}, we saw that the $\beta$ function of QED can
be computed from the counterterms for the electron-photon vertex, the
electron self-energy, and the photon self-energy.  The same derivation goes
through in the case of a non-Abelian gauge theory.  Thus, to lowest order,
\begin{equation}
  \beta(g) = gM\frac{\partial}{\partial M}\Bigl( -\delta_1 + \delta_2
    + \frac{1}{2}\delta_3 \Bigr) ,
  \label{eq:16.73}
\end{equation}
with the conventions for the counterterm vertices shown in
Fig.~\ref{fig:16.8}.  In QED, the first two terms cancel by the Ward
identity, so $\beta$ depends only on $\delta_3$.  In the non-Abelian case,
all three terms contribute.  The most difficult to compute is $\delta_3$, but
we have nearly done so already by computing the gauge-boson self-energy
diagrams.  Let us now complete this calculation of the $\beta$ function of
non-Abelian gauge theory.

\begin{figure}[t]
  \centering
  \includegraphics[width=0.6\textwidth]{fig16.8.png}
  \caption{Counterterms needed for computing fermion interactions in a
    non-Abelian gauge theory.}
  \label{fig:16.8}
\end{figure}

In order for the counterterm $\delta_3$ to cancel the divergence of
Eqs.~\eqref{eq:16.59} and \eqref{eq:16.71}, it must be of the form
\begin{equation}
  \delta_3 = \frac{g^2}{(4\pi)^2}\ \frac{\Gamma(2-d/2)}{(M^2)^{2-d/2}}\
    \Bigl( \frac{10}{3}C_2(G) - \frac{2}{3}n_fC(r) \Bigr) ,
  \label{eq:16.74}
\end{equation}
where $M$ is the renormalization scale.  Depending on the precise
renormalization conditions used, there may be additional finite contributions
to $\delta_3$, but these do not contribute to the $\beta$ function (to
one-loop order).  Similarly, the finite parts of $\delta_2$ and $\delta_1$
will depend on the details of the renormalization scheme.  However, as we
saw in Section~12.2, the one-loop contribution to the $\beta$ function is the
same in any scheme in which amplitudes are renormalized at a point where all
momentum invariants are of the same order $M^2$.  In dimensional
regularization, a logarithmic divergence always takes the form
$\Gamma(2-d/2)/\Delta^{2-d/2}$, where $\Delta$ is some combination of
momentum invariants.  Thus, to compute the $\beta$ function, we can simply
set $\Delta = M^2$ in such expressions.

\begin{figure}[t]
  \centering
  \includegraphics[width=0.55\textwidth]{fig16.9.png}
  \caption{Diagrams whose divergences are subtracted by the counterterms
    $\delta_2$ and $\delta_1$.}
  \label{fig:16.9}
\end{figure}

To complete the computation of the $\beta$ function, we must compute
$\delta_2$ and $\delta_1$ to the same level of approximation.  The fermion
self-energy counterterm $\delta_2$ cancels the divergence proportional to
$\slashed{k}$ in the first diagram of Fig.~\ref{fig:16.9}.  In Feynman-'t
Hooft gauge, the value of this diagram is
\begin{equation}
  \Sigma(k) = g^2C_2(r)\ \int_0^1 dx\ \int \frac{d^dp}{(2\pi)^d}\
    \frac{(d-2)\slashed{k}}{(p+k)^2p^2}\ \cdots
  = g^2C_2(r)\ \frac{\slashed{k}}{(4\pi)^2}\ \Gamma(2-d/2)\ \cdots .
  \label{eq:16.76}
\end{equation}
Since the divergence in the field strength renormalization is independent of
the fermion mass, we have simplified Eq.~(16.75) by setting the mass to
zero.  The product of group matrices equals the quadratic Casimir operator,
by definition \eqref{eq:15.92}.  The Dirac matrix structure can be reduced
using a contraction identity Eq.~(7.89).  The rest of the calculation
follows the same steps as for the boson self-energy diagrams.

The divergent part of this expression must be canceled by the second
counterterm diagram of Fig.~\ref{fig:16.8}.  Thus, if the renormalization
scale is $M$, the counterterm must be
\begin{equation}
  \delta_2 = \frac{g^2}{(4\pi)^2}\ \frac{\Gamma(2-d/2)}{(M^2)^{2-d/2}}\
    C_2(r) ,
  \label{eq:16.77}
\end{equation}
plus finite terms.  We note that, like $\delta_3$, $\delta_2$ depends on the
gauge; for example, $\delta_2$ has no one-loop divergence in Landau gauge
($\xi = 0$).

To determine $\delta_1$, we must compute the second and third diagrams of
Fig.~\ref{fig:16.9}.  The second diagram, computed in Feynman-'t Hooft gauge
and for massless fermions, is
\begin{equation}
  \Lambda^\mu = g^3\ \int \frac{d^dp}{(2\pi)^d}\ \frac{
    \gamma^\rho\slashed{k}\gamma^\sigma\ \gamma^\mu\ \gamma^\tau
    \slashed{k}\gamma^\upsilon }{(p+k')^2(p+k)^2p^2}\ \bigl[ C_2(r)
    - \tfrac{1}{2}C_2(G) \bigr]t^a\ \cdots ,
  \label{eq:16.78}
\end{equation}
where the gauge group matrices can be simplified according to
\begin{equation}
  \begin{aligned}
    t^bt^at^b &= t^bt^bt^a + t^b[t^a, t^b]\\
    &= C_2(r)t^a + it^bf^{abc}t^c\\
    &= C_2(r)t^a + \frac{1}{2}if^{abc}\cdot if^{bcd}t^d\\
    &= \bigl[ C_2(r) - \frac{1}{2}C_2(G) \bigr]t^a .
  \end{aligned}
  \label{eq:16.79}
\end{equation}
In the third line we have used the antisymmetry of $f^{abc}$ to rewrite the
matrix product as a commutator; in the last line we have used
Eq.~\eqref{eq:15.93}.

The diagrams computed earlier in this section had positive superficial
degrees of divergence, so we needed to extract their logarithmic divergences
carefully.  The integral in \eqref{eq:16.78}, however, is superficially
logarithmically divergent, and so the coefficient of this divergence can be
extracted easily by considering the limit in which the integration variable
$p$ is much greater than any external momentum.  In this limit, the diagram
is estimated as follows:
\begin{equation}
  \Lambda^\mu \approx g^3\bigl[ C_2(r) - \tfrac{1}{2}C_2(G) \bigr]t^a\ \int
    \frac{d^dp}{(2\pi)^d}\ \frac{\gamma^\rho\slashed{k}\gamma^\sigma\,
      \gamma^\mu\,\gamma^\tau\slashed{k}\gamma^\upsilon}{(p^2)^3}\ \cdots .
  \label{eq:16.80}
\end{equation}
If we replace $p^\rho p^\sigma$ by $g^{\rho\sigma}p^2/d$ in the numerator of
\eqref{eq:16.80}, this expression gives
\begin{equation}
  \Lambda^\mu = g^3\bigl[ C_2(r) - \tfrac{1}{2}C_2(G) \bigr]t^a\ \gamma^\mu\,
    \frac{3}{(4\pi)^2}\ \Gamma(2-d/2)\ \cdots .
  \label{eq:16.81}
\end{equation}
This estimate gives the correct coefficient of the divergent term.  It drops
completely the finite terms in the vertex function, but we do not need these
to compute the $\beta$ function.

The third diagram of Fig.~\ref{fig:16.9} can be analyzed in the same way.
Its value, in Feynman-'t Hooft gauge and for massless fermions, is
\begin{equation}
  \Lambda^\mu = g^3\ \int \frac{d^dp}{(2\pi)^d}\ \frac{\gamma^\rho
    \slashed{k}\gamma^\sigma\ \gamma^\mu }{(k'-p)^2(k-p)^2p^2}\ \cdot\
    \frac{1}{2}C_2(G)t^a\ \cdots ,
  \label{eq:16.82}
\end{equation}
where the gauge matrix product can be reduced as follows:
\begin{equation}
  f^{abc}t^bt^c = if^{abc}\cdot if^{bcd}t^d = C_2(G)t^a .
\end{equation}
Again we can determine the logarithmic divergence of this diagram by
neglecting all external momenta in comparison with $p$.  A straightforward
calculation then yields the divergent part.  Combining the divergent parts of
the three counterterms, we find
\begin{equation}
  \delta_1 = \delta_2 + \frac{1}{2}\delta_3\ \cdots ,
  \label{eq:16.89}
\end{equation}
up to gauge-dependent and finite terms.  Substituting into
\eqref{eq:16.73}, we obtain the famous result for the one-loop $\beta$
function of a non-Abelian gauge theory:
\begin{equation}
  \beta(g) = -\frac{g^3}{16\pi^2}\ \Bigl( \frac{11}{3}C_2(G)
    - \frac{4}{3}n_fC(r) \Bigr) .
  \label{eq:16.90}
\end{equation}

In our discussion of the counterterms of QED at the end of Section~10.3, we
remarked that the relation between $\delta_1$ and $\delta_2$ insured that all
electrically charged species see a common universal value of the coupling
constant $e$.  In non-Abelian gauge theories, the relations \eqref{eq:16.89}
and their higher-loop generalizations preserve the universality of the
non-Abelian couplings.  In QED, we were able to obtain an even stronger
relation, $\delta_1 = \delta_2$ or $Z_1 = Z_2$, from the absolute
normalization of the matrix elements of the vector current.  However, in
non-Abelian gauge theories, the corresponding vector current
$j^{\mu a} = \bar{\psi}\gamma^\mu t^a\psi$ transforms under local gauge
transformations in the adjoint representation.  Thus the Faddeev-Popov
prescription cannot be used to compute matrix elements of this current
unambiguously, and thus the normalization of these matrix elements is not
preserved by the perturbation theory.

\section{Asymptotic Freedom: The Background Field Method}

In the previous section, we saw that the $\beta$ function of a non-Abelian
gauge theory with a sufficiently small number of fermions is negative.  This
result is important enough that it is worthwhile to derive it twice.  The
preceding derivation was straightforward but not very illuminating.  In this
section we give a second derivation of the same result, which is more
abstract but much cleaner and more transparent.

The method of this section reflects the spirit of Wilson's idea of
integrating out the high-momentum degrees of freedom, while taking proper
care to preserve gauge invariance.  We will compute the effective action of a
non-Abelian gauge theory for a fixed, slowly varying, classical background
gauge field $\bar{A}^a_\mu(x)$.  By adopting a canonical normalization of this
field, we can interpret the coefficient of the effective action as a running
coupling constant.  This method is analogous to Polyakov's method for
computing the $\beta$ function of the nonlinear sigma model, presented in
Section~13.3.

\subsection{Background Field Perturbation Theory}

To set up the computation, rescale the gauge field $gA^a_\mu \to A^a_\mu$.  In
this normalization, the gauge coupling is removed from the covariant
derivative and moved to the coefficient of the gauge field kinetic energy
term.  We thus start from the Lagrangian
\begin{equation}
  \mathcal{L} = -\frac{1}{4g^2}(F^a_{\mu\nu})^2 + \bar{\psi}(i\slashed{D})\psi
    + \bar{c}^a(-\partial^\mu D_\mu)^{ab}c^b ,
  \label{eq:16.90}
\end{equation}
with
\begin{equation}
  D_\mu = \partial_\mu - iA^a_\mu t^a,\qquad
  F^a_{\mu\nu} = \partial_\mu A^a_\nu - \partial_\nu A^a_\mu
    + f^{abc}A^b_\mu A^c_\nu ,
  \label{eq:16.91}
\end{equation}
and the fermion mass set to zero for simplicity.  The transformation laws of
$A^a_\mu$ and $\psi$ are also independent of the coupling constant:
\begin{equation}
  \delta A^a_\mu = \partial_\mu\alpha^a + f^{abc}A^b_\mu\alpha^c,\qquad
  \delta\psi = i\alpha^a t^a\psi .
  \label{eq:16.92}
\end{equation}
On the other hand, the coupling constant $g$ will appear in the gauge field
propagator.

Next, split the gauge field into a classical background field and a
fluctuating quantum field:
\begin{equation}
  A^a_\mu \to \bar{A}^a_\mu + A^a_\mu .
  \label{eq:16.93}
\end{equation}
We will treat the classical part $\bar{A}^a_\mu$ as a fixed field
configuration and the fluctuating part $A^a_\mu$ as the integration variable
of the functional integral.  From here on, we will use the symbol $D_\mu$ to
denote the covariant derivative with respect to the background field:
$D_\mu = \partial_\mu - i\bar{A}^a_\mu t^a$.  Then
\begin{equation}
  \bar{\psi}(i\slashed{D})\psi \to \bar{\psi}(i\slashed{\partial})\psi
    + \bar{\psi}\slashed{A}^a t^a\psi .
  \label{eq:16.94}
\end{equation}
The Yang-Mills field strength decomposes as follows:
\begin{equation}
  \begin{aligned}
    F^a_{\mu\nu} &\to \partial_\mu\bar{A}^a_\nu - \partial_\nu\bar{A}^a_\mu
      + f^{abc}\bar{A}^b_\mu\bar{A}^c_\nu\\
    &\quad + \partial_\mu A^a_\nu - \partial_\nu A^a_\mu + f^{abc}\bigl(
      \bar{A}^b_\mu A^c_\nu - A^b_\nu\bar{A}^c_\mu \bigr) + f^{abc}A^b_\mu A^c_\nu\\
    &= \bar{F}^a_{\mu\nu} + D_\mu A^a_\nu - D_\nu A^a_\mu + f^{abc}A^b_\mu
      A^c_\nu ,
  \end{aligned}
  \label{eq:16.95}
\end{equation}
where, in the last line, $\bar{F}^a_{\mu\nu}$ is the field strength of the
classical field, and $D_\mu$ is the covariant derivative in the adjoint
representation, Eq.~\eqref{eq:15.86}.  Notice that, both in \eqref{eq:16.94}
and in \eqref{eq:16.95}, the derivative $\partial_\mu$ appears only as a part
of the covariant derivative with respect to the background field.

If the background field $\bar{A}^a_\mu$ is regarded as fixed, the Lagrangian
has a local gauge symmetry implemented by transformations on $A^a_\mu$:
\begin{equation}
  A^a_\mu \to A^a_\mu + D_\mu\alpha^a + f^{abc}A^b_\mu\alpha^c .
  \label{eq:16.96}
\end{equation}
To define the functional integral, we must gauge-fix using the Faddeev-Popov
procedure.  We choose a gauge-fixing condition that is covariant with respect
to the background gauge field:
\begin{equation}
  G^a(A) = D^\mu A^a_\mu - \omega^a ,
  \label{eq:16.97}
\end{equation}
instead of \eqref{eq:16.28}.  The Faddeev-Popov determinant involves the
variation of this operator with respect to the gauge transformation
\eqref{eq:16.96}.  As in Section~16.2, we can promote the gauge-fixing term
to the exponent, to quantize the theory in the background field analogue of
Feynman-'t Hooft gauge.  Then the gauge-fixed Lagrangian is
\begin{equation}
  \begin{aligned}
    \mathcal{L} &= -\frac{1}{4g^2}\bigl( \bar{F}^a_{\mu\nu} + D_\mu A^a_\nu
      - D_\nu A^a_\mu + f^{abc}A^b_\mu A^c_\nu \bigr)^2
      - \frac{1}{2}(D^\mu A^a_\mu)^2\\
    &\quad + \bar{\psi}(i\slashed{\partial} + \slashed{A}^a t^a)\psi
      - \bar{c}^a(D^2)^{ac}c^c - \bar{c}^a D^\mu\bigl( f^{abc}A^b_\mu c^c
        \bigr) .
  \end{aligned}
  \label{eq:16.98}
\end{equation}
The Lagrangian \eqref{eq:16.98} is gauge-fixed, but it is invariant under a
local symmetry that transforms both $A^a_\mu$ and the background field
$\bar{A}^a$:
\begin{equation}
  \begin{aligned}
    A^a_\mu &\to A^a_\mu - f^{abc}\beta^b A^c_\mu,\\
    \bar{A}^a_\mu &\to \bar{A}^a_\mu + D_\mu\beta^a,\\
    \psi &\to \psi + i\beta^a t^a\psi,\\
    c^a &\to c^a - f^{abc}\beta^b c^c .
  \end{aligned}
  \label{eq:16.99}
\end{equation}
Under this transformation, $A^a_\mu$ transforms as a matter field in the
adjoint representation, while $\bar{A}^a_\mu$ carries the part of the local
gauge transformation proportional to $\partial_\mu\beta^a$.  To prove that
\eqref{eq:16.99} is a symmetry of \eqref{eq:16.98}, we need only note that
\eqref{eq:16.98} is globally invariant, and that $\bar{A}^a_\mu$ appears in
\eqref{eq:16.98} only as a part of the covariant derivative and the field
strength.  The transformation \eqref{eq:16.99} is also a symmetry of the
functional measure.  Thus, if we functionally integrate over $A^a_\mu$,
$\psi$, and $c^a$ to compute the effective action, the result must be
invariant to local gauge transformations of $\bar{A}^a_\mu$.  This observation
greatly simplifies the analysis of the effective action.

\subsection{One-Loop Correction to the Effective Action}

Let us now compute the effective action, using the method of Section~11.4.
To compute $\Gamma[\bar{A}^a_\mu]$ to one-loop order, we drop terms linear in
the fluctuating field $A^a_\mu$ and then integrate over the terms quadratic in
$A^a_\mu$ and the fermion and ghost fields.  This produces functional
determinants, which we can evaluate into an appropriate form for an effective
action.

To carry out this program, we must work out the terms in \eqref{eq:16.98}
quadratic in each of the various fields.  The terms quadratic in $A^a_\mu$
are:
\begin{equation}
  \mathcal{L}_A = -\frac{1}{4}\Bigl( \frac{1}{2}(D_\mu A^a_\nu - D_\nu
    A^a_\mu)^2 + \bar{F}^a_{\mu\nu}f^{abc}A^b_\mu A^c_\nu + (D^\mu A^a_\mu)^2
    \Bigr) .
  \label{eq:16.100}
\end{equation}
After integrating by parts, we can rewrite this as
\begin{equation}
  \mathcal{L}_A = -\frac{1}{2g^2}\Bigl\{ A^a_\mu\bigl[ -(D^2)^{ab}g^{\mu\nu}
    + (D^\nu D^\mu)^{ab} - (D^\mu D^\nu)^{ab} \bigr]A^b_\nu
    - A^a_\mu f^{abc}\bar{F}^{b\mu\nu}A^c_\nu \Bigr\} .
  \label{eq:16.101}
\end{equation}
The term in brackets contains the commutator of covariant derivatives.  This
can be simplified using \eqref{eq:15.48}; the result combines with the last
term to give
\begin{equation}
  \mathcal{L}_A = -\frac{1}{2g^2}\Bigl\{ A^a_\mu\bigl[ -(D^2)^{ac}g^{\mu\nu}
    + 2f^{abc}\bar{F}^{b\mu\nu} \bigr]A^c_\nu \Bigr\} .
  \label{eq:16.102}
\end{equation}
The first term is part of a covariant d'Alembertian operator.  The second
term seems quite special, but we can put it into a form that will be
convenient later as follows: First, we recognize that $\bar{F}^{b\mu\nu}$ is
contracted with a group generator in the adjoint representation.  Next, we
introduce the matrix \eqref{eq:3.18} that is the generator of Lorentz
transformations on 4-vectors:
\begin{equation}
  (\mathcal{J}^{\rho\sigma})^\mu_{\ \nu} = i\bigl( \delta^{\rho\mu}
    \delta^{\sigma}_{\ \nu} - \delta^{\sigma\mu}\delta^{\rho}_{\ \nu}
    \bigr) .
  \label{eq:16.103}
\end{equation}
With these replacements, we can write \eqref{eq:16.102} in the form
\begin{equation}
  \mathcal{L}_A = -\frac{1}{2g^2}\Bigl\{ A^a_\mu\bigl[ -(D^2)^{ab}
    \delta^{\mu\nu} + 2\bigl( \tfrac{1}{2}\bar{F}^{b}_{\rho\sigma}
      \mathcal{J}^{\rho\sigma} \bigr)^{\mu\nu}(t^b_G)^{ac} \bigr]A^c_\nu
    \Bigr\} .
  \label{eq:16.104}
\end{equation}
The object in brackets can be considered as a generalized d'Alembertian for
fluctuations on the background field.

Next, we reduce the quadratic terms in fermion fields in a similar way.  The
quadratic Lagrangian for the fermion field is
\begin{equation}
  \mathcal{L}_\psi = \bar{\psi}(i\slashed{D})\psi .
  \label{eq:16.105}
\end{equation}
Integrating over the fermion fields, we find the determinant of the operator
$(i\slashed{D})$.  This is conveniently expressed as the square root of the
determinant of the operator
\begin{equation}
  \begin{aligned}
    (i\slashed{D})^2 &= \Bigl( -\frac{1}{2}\{ \gamma^\mu, \gamma^\nu \}
      - \frac{1}{2}[\gamma^\mu, \gamma^\nu] \Bigr)D_\mu D_\nu\\
    &= -D^2 + 2i\Bigl( -\frac{1}{2}[\gamma^\mu, \gamma^\nu] \Bigr)D_\mu D_\nu .
  \end{aligned}
  \label{eq:16.106}
\end{equation}
In the last line, the commutator of Dirac matrices forms the generator of
Lorentz transformations in the spinor representation, $S^{\mu\nu}$
\eqref{eq:3.23}.  Since this object is antisymmetric in its indices, the
product $D_\mu D_\nu$ that is contracted with it can be replaced by half of
their commutator.  Then \eqref{eq:16.106} takes the form
\begin{equation}
  (i\slashed{D})^2 = -D^2 + 2\bigl( \tfrac{1}{2}\bar{F}^{b}_{\rho\sigma}
    S^{\rho\sigma} \bigr)t^b ,
  \label{eq:16.107}
\end{equation}
where $t^a$ is now given in the representation of the fermions.  This is just
the d'Alembertian in \eqref{eq:16.104}, rewritten for the new set of spin and
gauge quantum numbers.  If the theory contains $n_f$ species of fermions, the
fermionic functional integral gives the determinant of \eqref{eq:16.107}
raised to the power $n_f/2$.

The quadratic term in ghosts is simply
\begin{equation}
  \mathcal{L}_c = \bar{c}^a\bigl[ -(D^2)^{ab} \bigr]c^b ,
  \label{eq:16.108}
\end{equation}
This contains the same d'Alembertian operator written for the case of spin
zero.

To summarize all of these results, we define the general covariant
background-field d'Alembertian as
\begin{equation}
  \Delta_{r,j} = -D^2 + 2\bigl( \tfrac{1}{2}F^b_{\rho\sigma}
    \mathcal{J}^{\rho\sigma}_j \bigr)t^b ,
  \label{eq:16.109}
\end{equation}
acting on a field of representation $r$ and spin $j$.  The square of the
covariant derivative gives the normal, convective, minimal coupling of the
particle described by $\Delta_{r,j}$ to the gauge field.  The additional term
is a magnetic moment interaction with the gauge field, whose strength
corresponds to a $g$-factor $g = 2$.  Using this general expression, we can
write the effective action for the classical fields $\bar{A}^a_\mu$ to
one-loop order, as
\begin{equation}
  e^{i\Gamma[\bar{A}]} = \det(\Delta_{G,1})^{1/2}\ \cdot\
    \bigl( \det\Delta_{r,1/2} \bigr)^{n_f/2}\ \cdot\
    \bigl( \det\Delta_{G,0} \bigr)^{-1}\ \cdot\ e^{i\int d^4x\
      \mathcal{L}_{c.t.}} ,
  \label{eq:16.110}
\end{equation}
where $\mathcal{L}_{c.t.}$ is the counterterm Lagrangian and the three
determinants are the results of evaluating the gauge field, fermion, and
ghost functional integrals.  Additional loop corrections to the effective
action are suppressed by another factor of $g^2$.

Since each integral contributing to \eqref{eq:16.110} is invariant to
\eqref{eq:16.99}, each determinant will be a gauge-invariant functional of
$\bar{A}^a_\mu$.  If we expand the determinants in powers of the background
field, we should then find a series of terms that begins
\begin{equation}
  \log\det\Delta_{r,j} = i\int d^4x\ \Bigl( \frac{1}{4}C_{r,j}(F^a_{\mu\nu})^2
    + \cdots \Bigr) ,
  \label{eq:16.111}
\end{equation}
where the succeeding terms contain higher-dimension gauge-invariant
operators.  The coefficient $C_{r,j}$ can depend on the representation $r$ and
the spin $j$.  This first term of the expansion modifies the zeroth-order
effective action according to
\begin{equation}
  \frac{1}{4g^2} \to \frac{1}{4g^2} + C_{r,j}\ \cdots .
  \label{eq:16.117}
\end{equation}
The factors $C_{r,j}$ are dimensionless but, since they arise from a
one-loop computation, we should expect that they are logarithmically
divergent:
\begin{equation}
  C_{r,j} = \bar{C}_{r,j}\ \log\frac{\Lambda^2}{M^2} + \cdots .
\end{equation}

\begin{figure}[t]
  \centering
  \includegraphics[width=0.55\textwidth]{fig16.10.png}
  \caption{Terms quadratic in the external field in the expansion of
    $\log\det\Delta_{r,j}$.  The special vertex arises from the
    $F^{\rho\sigma}\mathcal{J}_{\rho\sigma}$ coupling.}
  \label{fig:16.10}
\end{figure}

The terms in \eqref{eq:16.111} quadratic in $\bar{A}^a_\mu$ can be written in
Fourier space as
\begin{equation}
  \log\det\Delta_{r,j} = \frac{1}{2}\int \frac{d^4k}{(2\pi)^4}\
    \bar{A}^a_\mu(-k)\bar{A}^b_\nu(k)\ \bigl( k^2g^{\mu\nu} - k^\mu k^\nu
      \bigr)\ \cdot\ \bigl[ C_{r,j}\ \cdots \bigr] .
  \label{eq:16.121}
\end{equation}
We will now compute these terms explicitly from Eq.~(16.120) and bring
them into the form of \eqref{eq:16.121}.  The terms with two powers of
$A^a_\mu$ in the expansion Eq.~(16.120) are those with one power of
$\bar{A}^{(2)}$ or two powers of $\bar{A}^{(1)}$ or $\bar{A}^{(0)}$.  Further,
terms linear in $\bar{A}^{(0)}$ are proportional to $\mathrm{tr}[\mathcal{J}^{\rho\sigma}] = 0$, so the cross term between these two structures vanishes.  The three remaining contributions correspond to the Feynman diagrams shown in Fig.~\ref{fig:16.10}.

The term involving two powers of $\bar{A}^{(1)}$ is
\begin{equation}
  \frac{1}{2}\mathrm{tr}\bigl[ (-\partial^2)^{-1}\bar{A}^{(1)}
    (-\partial^2)^{-1}\bar{A}^{(1)} \bigr] = \cdots ,
  \label{eq:16.122}
\end{equation}
where the trace is now simply a trace over gauge and spin indices.  The
factor $1/2$ comes from the expansion of the logarithm.  The term involving
one power of $\bar{A}^{(2)}$ is
\begin{equation}
  \mathrm{tr}\bigl[ (-\partial^2)^{-1}\bar{A}^{(2)} \bigr] = \cdots .
  \label{eq:16.123}
\end{equation}
As Fig.~\ref{fig:16.10} suggests, these two contributions are precisely
proportional to the contribution of a scalar particle to the QED vacuum
polarization, times the factor
\begin{equation}
  \mathrm{tr}[t^at^b] = C(r)d(j)\delta^{ab} ,
  \label{eq:16.124}
\end{equation}
where $d(j)$ is the number of spin components.  The values of the diagrams
can be worked out using the methods of the previous section (or simply
recalled from Problem~9.1).  One finds that the two diagrams together sum up
to the gauge-invariant form \eqref{eq:16.121}, to give
\begin{equation}
  \frac{1}{2}\int \frac{d^4k}{(2\pi)^4}\ \bar{A}^a_\mu(-k)\bar{A}^b_\nu(k)\
    \bigl( k^2g^{\mu\nu} - k^\mu k^\nu \bigr)\ \cdot\
    \Bigl[ \frac{1}{3}d(j)C(r)\ \Gamma(2-d/2) + \cdots \Bigr] .
  \label{eq:16.125}
\end{equation}

The term involving two powers of $\bar{A}^{(0)}$ is
\begin{equation}
  -\frac{1}{2}\mathrm{tr}\bigl[ (-\partial^2)^{-1}\bar{A}^{(0)}
    (-\partial^2)^{-1}\bar{A}^{(0)} \bigr] = \cdots .
  \label{eq:16.126}
\end{equation}
To evaluate this, define $C(j)$ as the trace over spin indices
\begin{equation}
  \mathrm{tr}\bigl[ \mathcal{J}^{\rho\sigma}\mathcal{J}^{\alpha\beta} \bigr]
  = \bigl( g^{\rho\alpha}g^{\sigma\beta} - g^{\rho\beta}g^{\sigma\alpha}
    \bigr)C(j) .
  \label{eq:16.127}
\end{equation}
It is straightforward to work out from the explicit expressions that
\begin{equation}
  C(j) = \begin{cases}
    0 & \text{scalars};\\
    1 & \text{Dirac spinors};\\
    2 & \text{4-vectors}.
  \end{cases}
  \label{eq:16.128}
\end{equation}
Then \eqref{eq:16.126} can be evaluated as
\begin{equation}
  \frac{1}{2}\int \frac{d^4k}{(2\pi)^4}\ \bar{A}^a_\mu(-k)\bar{A}^b_\nu(k)\
    \bigl( k^2g^{\mu\nu} - k^\mu k^\nu \bigr)\ \cdot\
    \Bigl[ -4C(j)C(r)\ \Gamma(2-d/2) + \cdots \Bigr] .
  \label{eq:16.129}
\end{equation}
Adding \eqref{eq:16.125} and \eqref{eq:16.129}, we find that the coefficient
$C_{r,j}$ in \eqref{eq:16.111} is given by
\begin{equation}
  C_{r,j} = \Bigl( \frac{1}{3}d(j) - 4C(j) \Bigr)C(r)\ \Gamma(2-d/2) .
  \label{eq:16.131}
\end{equation}
Thus,
\begin{equation}
  C_{r,j} = \begin{cases}
    +\frac{1}{3}C(r) & \text{scalars};\\
    -\frac{8}{3}C(r) & \text{Dirac spinors};\\
    -\frac{20}{3}C(r) & \text{4-vectors}.
  \end{cases}
  \label{eq:16.132}
\end{equation}
Notice that, whenever the magnetic moment term is nonzero, it dominates, and
that its coefficient is opposite in sign from the convective term.

Inserting the values from \eqref{eq:16.132} into \eqref{eq:16.117}, we find
\begin{equation}
  \beta(g) = -\frac{g^3}{16\pi^2}\ \Bigl( \frac{11}{3}C_2(G)
    - \frac{4}{3}n_fC(r) \Bigr) .
  \label{eq:16.133}
\end{equation}
We thus confirm the conclusion of the previous section, that non-Abelian
gauge theories with sufficiently few fermions are asymptotically free.

\section{Asymptotic Freedom: A Qualitative Explanation}

In the previous two sections%
\footnote{Section~16.7 draws on the main result of 16.5 and 16.6, but does
not depend on these earlier sections.  However, even if you have not read
Section~16.5, you may wish to skim pages 522 through 531 to get an overview
of how the $\beta$ function can be calculated.}, we twice calculated the
$\beta$ function in non-Abelian gauge theory:
\begin{equation}
  \beta(g) = -\frac{g^3}{16\pi^2}\ \Bigl( \frac{11}{3}C_2(G)
    - \frac{4}{3}n_fC(r) \Bigr) .
  \label{eq:16.134}
\end{equation}
Here $n_f$ is the number of fermion species in representation $r$, $C(r)$ is
the constant appearing in the orthogonality relation \eqref{eq:15.78} for the
representation matrices, and $C_2(G)$ is the quadratic Casimir operator of
the adjoint representation of the group, defined in Eq.~\eqref{eq:15.92}.  In
an $SU(N)$ gauge theory with fermions in the fundamental representation, this
result becomes
\begin{equation}
  \beta(g) = -\frac{g^3}{16\pi^2}\ \Bigl( \frac{11}{3}N
    - \frac{2}{3}n_f \Bigr) .
  \label{eq:16.135}
\end{equation}
The overall minus sign implies that, for sufficiently small $n_f$, non-Abelian
gauge theories are asymptotically free.  In this case the running coupling
constant tends to zero at large momenta, according to Eq.~(12.92):
\begin{equation}
  \bar{g}^2(k) = \frac{\bar{g}^2(M)}{1 +
    \frac{\bar{g}^2(M)}{8\pi^2}\bigl( \frac{11}{3}N - \frac{2}{3}n_f
      \bigr)\log(k^2/M^2)} .
  \label{eq:16.136}
\end{equation}

The asymptotic freedom of non-Abelian gauge theories is a surprising
conclusion.  When we first encountered the running of the electromagnetic
coupling in Section~7.5, we found it easy to understand the direction of the
coupling constant flow: The vacuum acquires a dielectric property due to
virtual electron-positron pair creation, causing the effective electric charge
to decrease at large distances.  In non-Abelian gauge theories, according to
Eq.~\eqref{eq:16.134}, the fermions still produce such an effect.
Furthermore, since the non-Abelian gauge bosons are charged, they should
produce an additional screening effect.  But according to Eq.~\eqref{eq:16.134},
the net effect of the non-Abelian gauge bosons is opposite in sign.
Apparently there must be other, competing, effects, which overcome the effect
of screening and change the sign of the $\beta$ function.

The precise form of these effects depends on the gauge.  They are simplest to
describe in the Coulomb gauge, for which the gauge fixing condition is
\begin{equation}
  \partial^i A^{ai} = 0 .
  \label{eq:16.137}
\end{equation}
We will not work out the full quantization in this gauge; rather, we will
just describe its qualitative features.  As in electrodynamics, the field
quanta in Coulomb gauge are described in a non-Lorentz-covariant manner as
transversely polarized photons.  There are no timelike or longitudinal
photons and no propagating ghosts.  However, there is a Coulomb potential,
described by the field $A^{a0}$, which obeys a constraint equation analogous
to Gauss's law.  Not surprisingly, in the non-Abelian case, Gauss's law takes
the gauge-covariant form
\begin{equation}
  D^iE^{ai} = g\rho^a ,
  \label{eq:16.138}
\end{equation}
where $E^{ai} = F^{a0i}$ and $\rho^a$ is the charge density of the global
symmetry current of the fermions.  Recall from Eq.~\eqref{eq:15.86} that the
covariant derivative acting on a field in the adjoint representation is
\begin{equation}
  (D^i\phi)^a = \partial^i\phi^a + gf^{abc}A^{bi}\phi^c .
\end{equation}

To analyze the consequences of Eq.~\eqref{eq:16.138}, we will choose an
example as simple and explicit as possible.  Let the gauge group be $SU(2)$,
so that $a = 1, 2, 3$ and $f^{abc} = \epsilon^{abc}$.  Let us compute the
Coulomb potential of a point charge of magnitude $+1$ with the orientation
$a = 1$.  We will solve for $E^{ai}$ using an iteration process, putting the
gauge-field term of the covariant derivative on the right-hand side of the
equation:
\begin{equation}
  \partial^iE^{ai} = g\delta^{(3)}(\mathbf{x})\delta^{a1}
    + g\epsilon^{abc}A^{bi}E^{ci} .
  \label{eq:16.139}
\end{equation}
The second term on the right shows that, in a non-Abelian gauge theory, a
region containing vector potentials and electric fields that are parallel in
physical space and perpendicular in the group space is a source of electric
field.

\begin{figure}[t]
  \centering
  \includegraphics[width=0.55\textwidth]{fig16.11.png}
  \caption{The effect of vacuum fluctuations on the Coulomb field of an
    $SU(2)$ gauge theory.  In (a), a fluctuation $A^2$ occurs on top of the
    $1/r^2$ field $E^1$.  These combined fields generate a sink of the field
    $E^3$, as shown in (b).  The $E^3$ field, in turn, combines with $A^2$ to
    create an effective $E^1$ dipole, shown in (c).  The dipole points toward
    the original charge, enhancing its field at large distances.}
  \label{fig:16.11}
\end{figure}

The implication of Eq.~\eqref{eq:16.139} is worked out pictorially in
Fig.~\ref{fig:16.11}.  The leading term on the right-hand side of
\eqref{eq:16.139} implies a $1/r^2$ electric field of type $a = 1$ radiating
from $\mathbf{x} = 0$.  Somewhere in space, this electric field will cross
with a bit of vector potential $A^{ai}$ arising as a fluctuation of the
vacuum.  For definiteness, let us assume that this fluctuation has $a = 2$
and points in some diagonal direction, as shown in Fig.~\ref{fig:16.11}(a).
Then the second term on the right-hand side of Eq.~\eqref{eq:16.139} is
negative for $a = 3$: There is a sink of the field $E^{3i}$ at this location,
as shown in Fig.~\ref{fig:16.11}(b).  These new fields are, in two locations,
parallel or antiparallel to the original $A^{ai}$ field fluctuation.  Looking
again at the second term of Eq.~\eqref{eq:16.139}, we see that there is a
source of electric field with $a = 1$ closer to the origin, and a sink of
electric field with $a = 1$ farther away.  This is an induced electric dipole
in the vacuum, shown in Fig.~\ref{fig:16.11}(c).  But look at the signs: This
dipole is oriented toward the original charge, and thus serves to amplify
rather than screen it!  The effect of the original charge thus becomes
stronger at larger distances.

The competition between this antiscreening effect and the screening due to
virtual pairs of gauge bosons must be worked out quantitatively.  When this
is done,%
\footnote{T.~Appelquist, M.~Dine, and I.~Muzinich, \emph{Phys.\ Lett.}
\textbf{69B}, 231 (1977).} one finds that the antiscreening effect is 12
times larger.

In this argument, it is a set of dynamical features peculiar to the
non-Abelian gauge theory that enables the coupling constant to be amplified
rather than screened at large distances.  This suggests that asymptotic
freedom might be a special property of non-Abelian gauge theories.  Although
the statement can be proved only by exhausting other cases, it does actually
turn out to be true: Among renormalizable quantum field theories in four
spacetime dimensions, only the non-Abelian gauge theories are asymptotically
free.  We will use this property extensively in the next two chapters, where
we study Quantum Chromodynamics, the non-Abelian gauge theory of the strong
interactions.

\problems

\textbf{16.1} In this problem we will prove that the full Lagrangian
\eqref{eq:16.34} is invariant under the BRST transformation \eqref{eq:16.45}.
\begin{enumerate}[label=(\alph*)]
  \item Show that the transformation is nilpotent, that is, that applying it
    twice gives zero for each field.
  \item Show that the gauge-fixing term and the ghost term combine into a
    BRST-invariant combination.
\end{enumerate}

\textbf{16.2} In Section~16.5 we computed the one-loop $\beta$ function of a
non-Abelian gauge theory.  Verify the calculation of the gauge boson
self-energy by carrying out the integral in \eqref{eq:16.62} explicitly.

\textbf{16.3} \emph{The $\beta$ function in the background field method.}
Complete the derivation of the one-loop effective action in the background
field method, and verify that the coefficient $C_{r,j}$ in \eqref{eq:16.132}
leads to the $\beta$ function \eqref{eq:16.133}.

\textbf{16.4} \emph{One-loop divergences of the ghost sector.}  Compute the
one-loop ghost self-energy and the one-loop ghost-ghost-vector vertex, and
show that the divergences are canceled by the counterterms dictated by BRST
symmetry.

\textbf{16.5} In Section~16.7 we gave a qualitative argument for asymptotic
freedom in Coulomb gauge.  In this problem, work out the Coulomb gauge
quantization of the pure $SU(2)$ gauge theory, and compute the Coulomb
potential between two charges at large distances.
